%% n8a_eval.tex -- evaluation, first part: roadmap, design/measures, E1
%% (acquisition), E2 (organisation/lifecycle) and E3 (route-card audit).
%% Numbering follows the 2026-11-13 map (E3 = route audit, E4 = experience);
%% labels are historical and remain unique: sec:eval, sec:eval-setup, sec:e1,
%% sec:e2, sec:routing-audit, fig:routing, tab:eval-design.
\section{Evaluation}
\label{sec:eval}
We evaluate HMRS as a knowledge-management system with an access path, not as
an isolated ranker. Three evidence families are kept separate and are never
pooled into a system-level accuracy: management (E1--E2 and E4), access (E3
and E5) and answer/deployment (E6--E7), with E8 diagnosing verifier
placement and E9 bounding stability, concurrency, gate and scope behaviour.
E3 tests scripted collection choice and E5 measures retrieved text on the
same question set; E6 uses separate questions to inspect answers and traces.
In research-question terms, E1 and E2 answer RQ1, with E4 reporting the
operational experience store; E3 answers the collection-selection part of
RQ2, and E5's containment measure carries RQ2's document-narrowing clause;
E5 answers RQ3; E8 answers RQ4; E6 and E7
answer RQ5; E9 bounds all five questions. Table~S17 records the design and
provenance of each campaign; Tables~S9--S10 group the management
observations and retrieval designs with their limits.

\subsection{Experimental design and measures}
\label{sec:eval-setup}

\textbf{Corpora and questions.} Five data packages are used, each with its own
unit. The industrial demonstration corpus contains 42 documents in five
libraries, including twenty near-duplicate telemetry reports, and 16
developer-authored questions with gold document fragments; Q16 shares a token
with its gold filename. The arXiv campaign submitted 100 papers across about
25 domains and produced 322 stored parts in five libraries at the 2026-10-13
snapshot; its ten gold questions are keyed by arXiv identifiers. A later
re-check enumerated 328 parts in the same five libraries and audited their
descriptions; it is a separate state, not a matched rerun. The cross-lingual
ingestion corpus has 400 Wikipedia pages per language (English, Chinese and
Japanese) and 40 extractive queries per language; its index self-check is not
a task-level retrieval benchmark. A separate routing corpus has 138 documents
in 879 parts across ten domains, 72 in-domain questions and four
out-of-domain probes per arm; two cross-placeholder probes are excluded. The
same 72 questions were rerun as a same-question follow-up whose corpus grew
to 906 parts; it is not an independent sample. The packaged trace tables
retain 76 question IDs per complete arm; four R1--R4
cross-placeholder/out-of-domain probes remain in the CSV. They are excluded
from the 72-question in-domain containment and paired tests, and are not
pooled into an abstention-success rate. Questions and gold annotations were
developer-authored, not sampled user queries.

\textbf{Measures and controls.} Document-level Hit@5, Recall@5 and MRR@5
measure the returned top five, with MRR zero if no gold appears in that
window \cite{craswell-msmarco-2021}. Scripted route accuracy checks the
labelled collection, whereas annotated target-string containment checks
whether the target substring occurs anywhere in the returned evidence pack;
a sibling-document occurrence may count, so this is not gold-document
retrieval, entailment or semantic answer correctness. Automated RAG evaluation frameworks \citep{es-ragas-2024,%
saad-falcon-ares-2024} assess the semantic quality of generated answers,
while cumulated-gain measures \citep{jarvelin-cumulated-gain-2002} evaluate
graded relevance of ranked results; neither is used to score semantic answer
quality here. The answer grids check whether the gold identifier appears in
the answer or trace, whether at least $\lceil3n/5\rceil$ annotated terms are
present (strict) or an accepted alias (alias-aware), and whether a
document-style citation marker is present. A marker does not establish source
resolution or entailment. Wilson intervals are conditional on the questions
and run. Exact McNemar tests are exploratory, two-sided and unadjusted; no
familywise inference is intended across the listed contrasts. Agent arms use
one model and one run per condition; retries after transport failures are
not independent repeats. The answer-layer campaigns used the platform chat
API (Claude Code engine); the four-question harness case compared Claude
Code with oh-my-pi pinned to the ustc/deepseek-flash model (November 2026).
Two corpus-integrity defects were
repaired before measurement, so their natural failure rate is unknown.
Table~S17 records these boundaries; Tables~S9--S10 map
management functions and campaigns to their checks.

\subsection{E1: Acquisition and parsing}
\label{sec:e1}
The 100-paper campaign recorded 322 parts in five libraries. At the October
snapshot 91 paper names appeared in the ingest state and 79 papers were
retrievable: twelve recorded and nine unrecorded papers were not retrievable,
and the records do not establish per-paper parse status or the stage causing
each failure. A later re-check enumerated 328 parts in the same five
libraries and found a content-derived description for every part (328/328,
at least 60 characters, with no placeholder names), a metadata result
separate from retrieval accuracy. The same re-check ran a filename-stem
short-query identifier probe: a query built from each PDF filename stem was
searched in result names for the arXiv ID. The ID appeared for all 100
records after forced reindexing, explicit indexing and a probe-tool
correction. This is an identifier-visibility diagnostic, not scoped title
retrieval: 72 records did not specify a library and 45 returned a chunk-copy
name first. It is not a matched before--after comparison with the October
snapshot. A separate ten-question vector check, scoped to the five
libraries, placed the gold document first in all ten questions and invoked
the judge with its real engine in all ten; the retained summary does not
preserve complete judge inputs, so survivor and evidence-pack counts are not
reported as a 10/10 result. The cross-lingual test ingested 400 pages per
language with membership 0.875 for English and 1.0 for Chinese and Japanese;
its post-ingest self-retrieval is an index check, not a task-level
benchmark. A separate 50-paper, 26-domain campaign parsed and verified all
papers, but 41 negative tag flags did not block ingestion and no other
quality-review outcomes were recorded.

\subsection{E2: Organisation, consistency, and management operations}
\label{sec:e2}
The consistency checks show that the three-source state is auditable and
repairable, not that every operation is correct. A disposable-library rename
left stale chunks in the vector index: 18 instead of nine, with duplicate
results; forced re-index restored nine. The 322-part campaign passed its
content check on all parts and its three-source checks on a 33-part sample.
A disposable three-level fixture supported parent-to-child moves, listing,
reading and judging. Of 16 lifecycle scenarios, create, move, delete and
forced re-index passed, while task-status reporting and the update check
failed; merge and split were not controlled. Near-duplicate review over the
322-part corpus at a 0.90 threshold reported 35 candidate groups
(0.90--0.975) with zero true duplicates and deleted nothing (Table~S9).
Tag checking passed 165 parts,
while orphan tags and one cross-library registry leak were observed. The
graph supplies navigation rather than evidentiary support: a document move
changed four tag edges (4,675 to 4,679) on a 658-node graph, which is a
navigation record, not evidence. These results motivate reconciliation
without establishing its completeness.

\subsection{E3: Paired route-card audit: the description-layer contrast}
\label{sec:routing-audit}
The collection tier is audited as a paired description-layer contrast: route
cards carrying derived descriptions (A0) against nominally name-only cards
(A3), on one fixed question set. A scripted route-card probe returned the
labelled domain for 72/72 in-domain questions in A0 and 37/72 in nominally
name-only A3; 35 routes were correct only in A0
(Figure~\ref{fig:routing}a; exact McNemar $p=5.82\times10^{-11}$).
Evidence-pack containment was 37/72 for A0, 33/72 for A3 and 72/72 for the
mechanical read-all arm (Figure~\ref{fig:routing}b; A0--A3 10:6,
$p=0.4545$; read-all--A0 35:0, $p=5.82\times10^{-11}$). The frozen traces
recorded 31.4 KB and 16.3 s per A0 question, 309 KB and 13.8 s per A0h
question, and 24.3 MB and 253.7 s per read-all question
(Figure~\ref{fig:routing}c); the read-all scan issued 64{,}584 document-read
calls in total, about 897 per question. These are tool-side text and
tool-time observations, not network traffic, model input or end-to-end
latency. Read-all is a mechanical string scan, not an answering agent.

The contrast is a configuration diagnostic, not an independent agent-selection
test. An unlogged A3 fallback could read descriptions; some questions share
titles with route-card descriptions; A3 can search outside the evaluation
scope; and 52 of the 72 questions were automatically expanded, with 47 from
three domains. The difference therefore does not isolate the effect of
descriptions. It does show that the route card and the returned evidence are
separate tiers: A0 finds the collection in every question while A3 does not,
and the evidence pack can still miss the string when the collection is
correct (A0 37/72 containment). The read-all arm bounds what a full scan can
find, at about 78 times the recorded text of A0h; it says nothing about
careful reading or answer quality. The content-layer configurations, in
which the tiered chunk-window arm A0h reaches 51--52/72 containment, are
reported in Section~\ref{sec:ladder-opt}; the 37/72 figure above is this
audit's own route-card evidence-pack containment.

\begin{figure}[!t]
\centering
\includegraphics[width=\columnwidth]{fig_opt_route}
\caption{Route-card diagnostic on 72 in-domain questions. (a) Gold-domain
selection: A0 72/72, A3 37/72; Wilson 95\% intervals; 35 A0-only routes, exact
McNemar $p=5.82\times10^{-11}$. (b) Exact annotated target-string containment in the returned pack: A0 37/72,
A3 33/72, mechanical read-all 72/72; A0--A3 10:6 ($p=0.4545$), read-all--A0
35:0 ($p=5.82\times10^{-11}$). (c) Frozen trace-recorded text per question
(MB, log scale) against containment; A0h 0.309 MB/72.2\%, read-all 24.3
MB/100\%, A0 0.031 MB/51.4\%. In panel (c), the A0f and PD markers are
shifted horizontally by 8 pt solely for visibility; the summary data retain
their unshifted coordinates. Read-all is a mechanical string scan, not an
answering agent; A3 has an unlogged fallback and may search outside the
evaluation scope; tool time excludes answer generation. These are routing,
substring and text-volume diagnostics, not answer quality or an isolated
description effect.}
\label{fig:routing}
\end{figure}
